\def\tituloconcepto{Las puertas cuánticas son matrices unitarias}
% Si el capitulo se compila solo (sin main.tex), esta cabecera abre el documento.
% Antes de % Si el capitulo se compila solo (sin main.tex), esta cabecera abre el documento.
% Antes de % Si el capitulo se compila solo (sin main.tex), esta cabecera abre el documento.
% Antes de \input{cabecera} cada capitulo define \tituloconcepto.
\ifdefined\enmain\else
  \documentclass[11pt,a4paper]{article}
  \input{preambulo}
  \begin{document}
  \title{\tituloconcepto}
  \author{\autor}
  \date{\today}
  \maketitle
\fi
 cada capitulo define \tituloconcepto.
\ifdefined\enmain\else
  \documentclass[11pt,a4paper]{article}
  % Preambulo compartido por todos los documentos. Compilar con XeLaTeX/LuaLaTeX
% (o Tectonic). Si se usa Overleaf, elegir el compilador XeLaTeX.
\usepackage[spanish,es-noshorthands]{babel}
\usepackage{amsmath,amssymb}
\usepackage{booktabs,array}
\usepackage[a4paper,margin=2.5cm]{geometry}
\usepackage{xcolor}
\usepackage{listings}
\usepackage[most]{tcolorbox}
\usepackage{quantikz}
\usepackage[hidelinks]{hyperref}

% ---- Autor (editar a gusto) -------------------------------------------------
\newcommand{\autor}{Genesis Silesky}

% ---- Notacion de Dirac y atajos ---------------------------------------------
% quantikz ya define \ket y \bra; DeclareRobustCommand los redefine sin error.
\DeclareRobustCommand{\ket}[1]{\left|#1\right\rangle}
\DeclareRobustCommand{\bra}[1]{\left\langle#1\right|}
\DeclareRobustCommand{\braket}[2]{\left\langle#1\middle|#2\right\rangle}
\newcommand{\C}{\mathbb{C}}
\newcommand{\R}{\mathbb{R}}
\DeclareMathOperator{\Tr}{Tr}

% ---- Cajas de apoyo --------------------------------------------------------
% "concepto": definiciones previas que hacen falta para entender el capitulo.
\newtcolorbox{concepto}[1]{breakable,colback=blue!4,colframe=blue!55!black,
  fonttitle=\bfseries,title={Concepto previo: #1}}
% "nota": aclaraciones y advertencias.
\newtcolorbox{nota}[1]{breakable,colback=orange!6,colframe=orange!70!black,
  fonttitle=\bfseries,title={#1}}

% ---- Estilos de listados ----------------------------------------------------
\lstdefinestyle{py}{language=Python,basicstyle=\ttfamily\footnotesize,
  breaklines=true,postbreak=\mbox{\textcolor{gray}{$\hookrightarrow$}\space},
  numbers=left,numberstyle=\tiny\color{gray},frame=single,
  commentstyle=\color{green!40!black}\itshape,
  keywordstyle=\color{blue!70!black}\bfseries,
  stringstyle=\color{red!60!black},showstringspaces=false,tabsize=4,
  columns=fullflexible,keepspaces=true}
\lstdefinestyle{salida}{basicstyle=\ttfamily\footnotesize,breaklines=true,
  frame=leftline,backgroundcolor=\color{gray!8},columns=fullflexible,
  keepspaces=true}

  \begin{document}
  \title{\tituloconcepto}
  \author{\autor}
  \date{\today}
  \maketitle
\fi
 cada capitulo define \tituloconcepto.
\ifdefined\enmain\else
  \documentclass[11pt,a4paper]{article}
  % Preambulo compartido por todos los documentos. Compilar con XeLaTeX/LuaLaTeX
% (o Tectonic). Si se usa Overleaf, elegir el compilador XeLaTeX.
\usepackage[spanish,es-noshorthands]{babel}
\usepackage{amsmath,amssymb}
\usepackage{booktabs,array}
\usepackage[a4paper,margin=2.5cm]{geometry}
\usepackage{xcolor}
\usepackage{listings}
\usepackage[most]{tcolorbox}
\usepackage{quantikz}
\usepackage[hidelinks]{hyperref}

% ---- Autor (editar a gusto) -------------------------------------------------
\newcommand{\autor}{Genesis Silesky}

% ---- Notacion de Dirac y atajos ---------------------------------------------
% quantikz ya define \ket y \bra; DeclareRobustCommand los redefine sin error.
\DeclareRobustCommand{\ket}[1]{\left|#1\right\rangle}
\DeclareRobustCommand{\bra}[1]{\left\langle#1\right|}
\DeclareRobustCommand{\braket}[2]{\left\langle#1\middle|#2\right\rangle}
\newcommand{\C}{\mathbb{C}}
\newcommand{\R}{\mathbb{R}}
\DeclareMathOperator{\Tr}{Tr}

% ---- Cajas de apoyo --------------------------------------------------------
% "concepto": definiciones previas que hacen falta para entender el capitulo.
\newtcolorbox{concepto}[1]{breakable,colback=blue!4,colframe=blue!55!black,
  fonttitle=\bfseries,title={Concepto previo: #1}}
% "nota": aclaraciones y advertencias.
\newtcolorbox{nota}[1]{breakable,colback=orange!6,colframe=orange!70!black,
  fonttitle=\bfseries,title={#1}}

% ---- Estilos de listados ----------------------------------------------------
\lstdefinestyle{py}{language=Python,basicstyle=\ttfamily\footnotesize,
  breaklines=true,postbreak=\mbox{\textcolor{gray}{$\hookrightarrow$}\space},
  numbers=left,numberstyle=\tiny\color{gray},frame=single,
  commentstyle=\color{green!40!black}\itshape,
  keywordstyle=\color{blue!70!black}\bfseries,
  stringstyle=\color{red!60!black},showstringspaces=false,tabsize=4,
  columns=fullflexible,keepspaces=true}
\lstdefinestyle{salida}{basicstyle=\ttfamily\footnotesize,breaklines=true,
  frame=leftline,backgroundcolor=\color{gray!8},columns=fullflexible,
  keepspaces=true}

  \begin{document}
  \title{\tituloconcepto}
  \author{\autor}
  \date{\today}
  \maketitle
\fi


\section{Las puertas cuánticas son matrices unitarias}

\subsection{Objetivo}
Ver que aplicar una puerta equivale a multiplicar una matriz por el vector de
estado, y por qué esa matriz debe ser unitaria.

\subsection{Conceptos previos}
\begin{concepto}{multiplicación matriz--vector y unitariedad}
Ver sección~\ref{sec:prerreq}: matrices unitarias ($U^\dagger U=I$) y notación
de Dirac. Recordatorio: el producto
$\begin{psmallmatrix}a&b\\c&d\end{psmallmatrix}\binom xy=\binom{ax+by}{cx+dy}$.
\end{concepto}

\subsection{Fundamento teórico}
Una puerta de un qubit es una matriz $U\in\C^{2\times2}$ unitaria y el estado
evoluciona como $\ket{\psi'}=U\ket\psi$. Las más usadas:
\begin{equation}
H=\frac1{\sqrt2}\begin{pmatrix}1&1\\1&-1\end{pmatrix},\quad
X=\begin{pmatrix}0&1\\1&0\end{pmatrix},\quad
Z=\begin{pmatrix}1&0\\0&-1\end{pmatrix}.
\end{equation}
Con $H\ket0=\frac1{\sqrt2}(\ket0+\ket1)=:\ket+$ y
$H\ket1=\frac1{\sqrt2}(\ket0-\ket1)=:\ket-$. \emph{(Aquí se usa la notación
$\ket\pm$ solo como abreviatura.)}

\paragraph{Por qué unitarias.} Si $U^\dagger U=I$ entonces
$\|U\ket\psi\|^2=\bra\psi U^\dagger U\ket\psi=\|\ket\psi\|^2$: la
normalización (probabilidades que suman 1) se conserva. Además $U$ es
invertible con $U^{-1}=U^\dagger$, es decir, la evolución cuántica es
\emph{reversible}. Para $H$ se cumple $H^\dagger=H$ y $H^2=I$.

\subsection{Código (completamente comentado)}
\lstinputlisting[style=py]{../../fuentes/02_puertas_matrices.py}

\subsection{Explicación celda por celda}
\begin{description}
  \item[Celda 1] \texttt{Operator} convierte un circuito en su matriz.
  \item[Celda 2] Para cada puerta crea un circuito de un qubit y extrae su
    matriz. \texttt{getattr(qc, "h")(0)} equivale a \texttt{qc.h(0)}; se usa
    para recorrer varias puertas en un ciclo.
  \item[Celda 3] Calcula $H\ket0$ con \texttt{@} (producto matricial de
    NumPy) y con \texttt{.evolve()} de Qiskit; ambos dan
    $(1/\sqrt2,\,1/\sqrt2)\approx(0.707,\,0.707)$.
  \item[Celda 4] Verifica numéricamente $H^\dagger H=I$ y $H^2=I$.
    \texttt{np.allclose} compara con tolerancia numérica (no se usa
    \texttt{==} con decimales).
\end{description}

\subsection{Salida obtenida}
\lstinputlisting[style=salida]{../../fuentes/salidas/02_puertas_matrices.txt}

\subsection{Interpretación}
Cada puerta es una matriz concreta y calculable. El resultado
$H\ket0=(0.707,0.707)$ confirma una superposición con
$P(0)=P(1)=\tfrac12$.

\subsection{Ejercicios}
\begin{enumerate}
  \item Calcule a mano $XZ\ket0$ y $ZX\ket0$. ¿Conmutan $X$ y $Z$?
  \item Verifique en el código que $X$ y $Z$ son unitarias.
  \item Calcule $H\ket1$ y compruebe que sus probabilidades son $\tfrac12,\tfrac12$
    aunque una amplitud sea negativa.
\end{enumerate}

\subsection{Errores comunes}
\begin{itemize}
  \item Con varios qubits, la matriz de un circuito usa el orden
    little-endian de Qiskit; para un solo qubit no hay ambigüedad.
  \item Comparar matrices con \texttt{==}: use \texttt{np.allclose}.
\end{itemize}

\ifdefined\enmain\else
  \end{document}
\fi

